%! Author = Sujal Singh
%! Date = 10/23/23

% Preamble
\documentclass[11pt,a4paper,oneside]{scrarticle}

% Packages
\usepackage{mathptmx}
\usepackage{tabularray}
\usepackage{graphicx}
\usepackage{sectsty}

\sectionfont{\centering}

\renewcommand*{\maketitle}{
    \pagenumbering{gobble}
    \begin{center}
        \vspace*{8em}
        \includegraphics[width=200pt]{../../tools/pdf/template/logo} \\[40pt]
        \textbf{\Huge LLB Assignment} \\[40pt]
        \begin{table}[htb]
            \label{tab:name-slip}
            \begin{tblr}{rows={30pt},colspec={|X|X[3]|},hlines}
                \textbf{Course}            & Industrial Internet of Things \\
                \textbf{Batch}             & 1                             \\
                \textbf{Name}              & Sujal Singh                   \\
                \textbf{Enrollment Number} & XXXXXXXXXXXX                  \\
            \end{tblr}
        \end{table}
    \end{center}
    \newpage
    \pagenumbering{arabic}
}

% Document
\begin{document}
    \maketitle

    \begin{center}
        \textbf{\LARGE Violations of Freedom of Speech and Expression in the Digital World by the Indian Government.}
        \\[10pt]
    \end{center}
    \begin{flushright}
        \textit{--- Sujal Singh} \\[20pt]
    \end{flushright}

    %------------------------------------------------------------------------------------------------------------------%
    %------------------------------------------------------------------------------------------------------------------%

    \section*{Abstract}\label{sec:abstract}
    The Internet is one of mankind's greatest achievements and an open Internet, free of censorship and control by any
    authority, is crucial to avoid a dystopian future, but ever since it's inception, government authorities have
    employed various tactics to censor the information published via the World Wide Web.\ This report details various
    acts by the Indian government supporting censorship on the Internet and using it as a tool for mass-surveillance,
    such acts go against constitutional values as envisioned by the authors of the Indian Constitution, and people in
    relevant fields of studies should feel responsible to oppose the continued decline of the open Internet, due to the
    government and massive corporations, and they should push for technical solutions that uphold constitutional values.
\end{document}